\section*{Annex I: Clues from Playtests}
\addcontentsline{toc}{section}{Annex I: Clues from Playtests}

The following catalogue is the output of the three-phase systematic analysis described in the Methods chapter: initial descriptive coding of field notes and debrief responses, comparative organization around continuity, difference, and novelty, and articulation of design principles connecting mechanics to observed effects.

Unlike a traditional results catalogue, which reports only what happened, this annex is organized around a pre-defined taxonomy of \textit{potential clues}: the observable signs that would count as evidence \textit{for} or \textit{against} the format's intended effects across each evaluation field. This taxonomy was specified \textit{before} data collection, as described in the Methods chapter, to ensure that analysis remained responsive to disconfirming evidence and did not converge only on confirming observations. For each potential clue, the catalogue records whether it was observed during the playtests, with what evidential support, and at what strength. Clues that did not manifest are recorded explicitly: the absence of a negative clue is evidence of the format functioning as intended; the absence of a positive clue is a genuine gap.

Evidence is classified by direction and strength. \textbf{Positive clues} are behaviors or reports consistent with the format achieving its intended effect. \textbf{Negative clues} are behaviors or reports indicating a design gap or failure. Strength is assessed as \textit{Strong} (behavioral, multi-source, or unsolicited), \textit{Moderate} (self-reported or single-source), or \textit{Weak} (speculative or highly context-dependent). Consistent with the comparative analysis described in the Methods chapter, behavioral and emergent evidence carries more weight than self-reported responses, and novel or different behaviors (departures from standard debate practice) constitute stronger design evidence than continuities, which may reflect participants' existing training rather than effects of the new format. Self-reported evidence gathered in response to direct questions in the semi-structured group debrief is noted where relevant, as it remains evidentially valid even if weaker than unsolicited or behavioral clues.

\textit{Sources: P1 = First playtest field notes (Jean Monnet debate club, March 2026); P2 = Second playtest field notes (April 2026); D1/D2 = Post-session debrief responses after first/second playtest; T2 = Transcript of second playtest.}

\bigskip

\noindent\textbf{A note on social desirability in Field 1.} Engagement and popularity clues are the most susceptible to social desirability effects: participants knew they were playtesting a format designed by the researcher in the room, and the exclusive nature of the invitation may have elevated their investment in displaying enthusiasm. Positive engagement evidence should therefore be read with additional skepticism; behavioral indicators are weighted higher than self-reports, and negative clues, which participants would have had more reason to suppress, carry elevated weight when they do appear.

\bigskip

%%% FIELD 1 %%%

\noindent\textit{Why Field 1 is a required characteristic.} Transformative potential cannot manifest in a game that players do not return to. The research question asks how game design theory can elevate a format's transformative potential, but "elevated transformative potential" presupposes that participants inhabit the format long enough for its mechanisms to produce their effects. The transformational container forms and deepens through repeated participation; bleed arcs, skill accumulation, and epistemic shifts are all longitudinal processes that require the format to generate genuine voluntary re-engagement. A format that produces one interesting session but fails to motivate return has demonstrated novelty, not transformative design. Field 1 therefore functions as the necessary entry condition for all subsequent fields: every transformative effect assessed in Fields 2--5 is predicated on the format being one people choose to play again.

\bigskip

\subsection*{Field 1: Engagement, Appeal, and Long-Term Viability}

\subsubsection*{Positive potential clues}

\begin{enumerate}

\item \textbf{Unsolicited expressions of enthusiasm during or after the debate}\\
\textit{Points toward:} Positive affective engagement; potential for voluntary repeated participation.\\
\textit{Relevance to field:} Unsolicited expressions isolate intrinsic motivation as the causal driver, removing the social-obligation confound that makes prompted enthusiasm ambiguous. For Field 1, intrinsic motivation is the operative mechanism: a format people are enthusiastically drawn toward without external encouragement is one they have reason to choose on their own, which is the basis of the voluntary repeated participation that long-term viability requires.\\
\textit{Observed:} Yes (P1, P2, D1, D2).\\
\textit{Evidence:} Participants described the experience as ``surprisingly liberating'' (P1, D1) and ``very very fun'' during the Huddle phase (P2, D2). Multiple unprompted positive comments recorded in both debrief forms without being solicited.\\
\textit{Mechanic:} Removal of binary win/loss outcomes, combined with out-of-sight Evaluator seating, reduces performance-oriented anxiety; personal dilemma motions create genuine stakes without threatening social standing, making organic enthusiasm more likely.\\
\textit{Strength:} Moderate. Self-reported; subject to social desirability effects; see caveat above.

\item \textbf{Passionate or animated discussion among participants without researcher prompting}\\
\textit{Points toward:} Genuine peer engagement; format generates internal conversational momentum.\\
\textit{Relevance to field:} Discussion that sustains itself beyond the formal debrief, without researcher steering, demonstrates that the format's content generates its own conversational momentum independent of structure. For Field 1, the key indicator is that engagement outlasts the session: a format that continues to be discussed after the researcher has stopped directing the interaction has generated the pull that makes planning future sessions feel like an obvious next step.\\
\textit{Observed:} Yes (P1, P2).\\
\textit{Evidence:} Post-debrief conversation in both sessions extended organically beyond the structured questions. In the second playtest, participants continued discussing the motion and their own values informally after the session ended. Both sessions also produced animated group conversations specifically about the experience of using second-person direct address and personal appeals during speeches: participants in both groups noted that this mode of speaking (addressing the Subject directly as ``you'' and appealing to their personal circumstances) felt novel, engaging, and surprisingly enjoyable to both perform and receive.\\
\textit{Mechanic:} Personal dilemma motions make participants' own values directly relevant, creating natural grounds for continued conversation; the debrief structure designates and legitimizes extended reflection rather than truncating it. The second-person address mode inverts standard debate's impersonal register, making the emotional texture of the debate itself a topic of enthusiastic reflection.\\
\textit{Strength:} Moderate. Observed; behavioral; but also consistent with social novelty of the event.

\item \textbf{Visible engagement signals during debate itself (laughter, focused silence, forward leaning)}\\
\textit{Points toward:} Format sustaining in-the-moment attention and flow.\\
\textit{Relevance to field:} In-the-moment engagement signals confirm that the experience of playing, rather than only the memory of having played, is positively absorbing. For Field 1, this is the real-time counterpart to post-session enthusiasm: a format that participants remember fondly but are visibly disconnected from during play has a mismatch between anticipated and actual experience that would degrade over repetition.\\
\textit{Observed:} Yes (P1 and P2).\\
\textit{Evidence:} Second playtest field notes record observable engagement during speeches and very novel and creative decisions: to play background music during S1 speech, to adopt speaking accents etc. In the first playtest, as players became more aware of what they could do in the debate, they started to have more fun and more passionate speeches.\\
\textit{Mechanic:} The Huddle phase builds emotional investment before speeches begin through interactive preparation; the Subject as a live, named decision-maker provides a concrete human target rather than an abstract adjudicator, sustaining speaker investment.\\
\textit{Strength:} Moderate. Observational; first-playtest engagement was inconsistent.

\item \textbf{Unprompted recognition of the format's benefits}\\
\textit{Points toward:} Participants forming their own positive evaluation independent of researcher framing.\\
\textit{Relevance to field:} When participants independently articulate why the format is valuable, they are constructing the self-sustaining narrative that motivates voluntary future participation and recruits new players. For Field 1, long-term viability at community scale depends on participants having reasons they can articulate to others, not only on individual desire to return.\\
\textit{Observed:} Yes (P1, P2, D1, D2).\\
\textit{Evidence:} Multiple participants articulated benefits spontaneously: accessibility, empathy development, skill transfer, reduced anxiety. In the second playtest, there was conversation on how the format would be a useful personal development tool for ``corpo trainings.''\\
\textit{Mechanic:} Not tied to one particular mechanic but rather the format as a whole.\\
\textit{Strength:} Moderate. Self-report; theoretically coherent; partially behavioral (the personal-development-tool framing implies intention to act).

\item \textbf{Expressed desire to replay or requests to organize future rounds}\\
\textit{Points toward:} Replayability; format generates sustained rather than one-off interest.\\
\textit{Relevance to field:} Expressed desire to replay is the most proximate forward-looking indicator of long-term viability available within a single-session playtest. For Field 1, this clue also captures whether participants perceive an ongoing growth arc within the format: naming what they would do differently next time signals that the format has made its own depth visible, which is the motivational basis for sustained multi-session engagement.\\
\textit{Observed:} Yes (P1, D1).\\
\textit{Evidence:} Participants said they would play again and would aim to be ``more emotionally immersed'' next time (P1, D1). In the second playtest players reported they would find it a fun game especially to play with friends, as a passtime, and to get their "build" done.\\
\textit{Mechanic:} The roleplay dimension of the format (emotional inhabitation of a character's values) is recognized by participants as an unexplored skill, creating an intrinsic learning arc; the commendation system provides named targets (e.g., ``Believable Passion'') for improvement across sessions.\\
\textit{Strength:} Moderate. Self-reported intention; actual re-engagement behavior untested.

\item \textbf{High enjoyment scores when directly asked}\\
\textit{Points toward:} Positive affective engagement; format generating a genuinely enjoyable experience rather than merely tolerable or dutiful participation.\\
\textit{Relevance to field:} Enjoyment is the affective substrate of voluntary repeated participation: people return to activities they found genuinely pleasant, and a format that is intellectually interesting but unpleasant to play will be studied rather than played. For Field 1, consistently high scores across participants indicate the format reliably produces positive experience, not only occasional highlights.\\
\textit{Observed:} Yes (D1, D2).\\
\textit{Evidence:} When asked directly in the debrief form how much they enjoyed the game, participants in both sessions reported very high scores. While solicited self-report is weaker evidence than unsolicited enthusiasm, the consistently high numbers across both groups, including participants who were not particularly enthused about other aspects, indicate that the enjoyment response was not merely social desirability but genuine. The convergence with unsolicited expressions in clue 1 above strengthens the overall picture.\\
\textit{Mechanic:} Same as clue 1.\\
\textit{Strength:} Moderate. Direct debrief form response; subject to social desirability effects; convergent with unsolicited expressions.

\end{enumerate}

\subsubsection*{Negative potential clues}

\begin{enumerate}

\item \textbf{Taciturn, flat, or deflated behavior after the round}\\
\textit{Points away from:} Engagement and emotional investment; would suggest the format failed to generate genuine stakes.\\
\textit{Relevance to field:} Post-round affect is the most immediate behavioral indicator of whether the format generated genuine investment rather than dutiful participation. For Field 1, deflated post-round behavior would signal the experience left no residue worth processing, no sense of unfinished business or growth arc to continue, making voluntary return unlikely.\\
\textit{Observed:} No.\\
\textit{Evidence:} Not manifested across either playtest. Post-round behavior in both sessions was animated and conversational.\\
\textit{Strength:} N/A.

\item \textbf{Direct criticism of the format itself}\\
\textit{Points away from:} Overall appeal; structural design coherence.\\
\textit{Relevance to field:} Structural criticism identifies design barriers to adoption and re-engagement. For Field 1, usability problems are long-term viability problems: a format with persistent confusing or frustrating elements will lose players between sessions even if the core experience motivates return, because the overhead of relearning and re-onboarding discourages casual re-engagement.\\
\textit{Observed:} Partially (P1, D1).\\
\textit{Evidence:} The commendation system was criticised as ``overwhelming'' and difficult to use, particularly the size of the badge list (P1, D1). This was a targeted structural criticism rather than a wholesale rejection. The issue of commendation size was not directly addressed by a design change; it is to some extent intrinsic to the system. However, participants recognized that cognitive load would decrease with familiarity: as an Evaluator learns the commendation list over repeated use, the overhead diminishes naturally. Participants also noted that a practical workaround is available from the start: an Evaluator can simply observe what debaters do well in their speeches and work backwards to identify which commendation matches, rather than scanning the list prospectively. This makes the system more accessible at small scale than it first appears.\\
\textit{Mechanic:} The commendation list in the first design iteration (Annex A) was designed for tournament scale, where Evaluators memorize a limited set of commendations over repeated use; at small informal scale with an unfamiliar system, the cognitive overhead of tracking all options exceeded what the context could support. The issue is thus primarily one of onboarding rather than structural incoherence.\\
\textit{Strength:} Moderate. Specific and actionable; confined to one subsystem; not generalized to the format overall; partially mitigated by intrinsic learning curve and retrospective matching strategy.

\item \textbf{Visible loss of focus or disengagement during speeches}\\
\textit{Points away from:} Sustained attention; flow state during play.\\
\textit{Relevance to field:} In-session disengagement indicates the format is failing to sustain attentional pull in real time. For Field 1, a format that participants remember fondly but were visibly disconnected from while playing has an experience-memory gap that narrows over repetition: disengagement during play would eventually erode even enthusiastic retrospective impressions.\\
\textit{Observed:} Partially (P1).\\
\textit{Evidence:} One participant appeared uncertain about what to do during the Subject role in the first playtest, producing visibly hesitant behavior. This was role-confusion-driven rather than disengagement with the format broadly, and was not repeated in the second playtest.\\
\textit{Mechanic:} The S0 and S1 instructions in the first design iteration did not specify the Subject's active task during other teams' speeches, leaving the Subject without a clear role during the middle of the debate.\\
\textit{Strength:} Weak. Single instance; role-specific; addressed in next design iteration.

\item \textbf{Expressed preference for returning to standard competitive debate}\\
\textit{Points away from:} Standalone appeal; would suggest format is only tolerable as an alternative, not as a preferred experience.\\
\textit{Relevance to field:} Expressed preference for existing formats is the most direct comparative appeal test: it measures whether Inside Out is perceived as genuinely distinct and preferable, or as a tolerable but inferior alternative. For Field 1, long-term viability in a community that already has established format options depends on the format offering something those options do not.\\
\textit{Observed:} Partially (P1, D1, P2, D2).\\
\textit{Evidence:} Across both playtests, a total of three participants (one in the first session and two in the second) expressed reservations tied to the format's non-competitive structure. All three were men who had been highly competitive and quite successful in traditional debate. In the first session, one participant described the format as ``anticlimactic'' and expressed missing the ``ego-boosting'' feeling of decisive victory (P1, D1); several others in that session suggested the format might be ``easier'' or ``for beginners,'' raising concern it could be perceived as less serious by the broader debate community. In the second session, in addition to similar sentiments, one participant raised a more structural concern: that the format's apparent benefits might be partly an artifact of playing without an established meta: as a meta develops around Inside Out (as it has in World Schools and BP), the same competitive pressures and strategic distortions that the format tries to eliminate could re-emerge. This is a theoretically significant observation worth tracking in future iterations.\\
\textit{Mechanic:} The removal of binary win/loss outcomes and zero-sum scoring is a deliberate design choice; these negative responses confirm that mechanic is operating as designed, not malfunctioning. The format is not intended to serve this participant profile as a primary audience, though the meta-concern raised in D2 is a genuine design question about long-term robustness.\\
\textit{Strength:} Moderate for scoping the target audience; weak as evidence of a format-wide problem. The meta-concern is Weak at this stage (speculative, without empirical basis) but theoretically interesting.

\item \textbf{``Cringe'' reactions to ritual elements or format conventions}\\
\textit{Points away from:} Immersive design; willingness to engage with the format's fictional premise.\\
\textit{Relevance to field:} Ritual failure makes the format feel awkward or embarrassing to inhabit, and willingness to perform awkward conventions is higher among enthusiasts than among casual players. For Field 1, persistent ritual cringe would progressively narrow the format's accessible audience, capping the long-term viability it needs to develop its transformative potential at scale.\\
\textit{Observed:} Yes (P1, D1, P2, D2).\\
\textit{Evidence:} The ending ritual (Rolling Out) was described by participants as ``very robotic'' and as something that ``takes you out of the immersion'' (P1, D1). This reaction recurred in the second playtest: participants were observed mumbling the scripted phrases rather than delivering them with investment, and the researcher (as designer) was asked by participants to justify the purpose of the ritual phrases, an indication of a breakdown in the fictional frame (P2, D2). This is the strongest negative engagement clue from either playtest. It indicates that immersion \textit{was} present before that point (participants could be ``in the story'') but was disrupted by the ritual structure. The ritual was redesigned in the third design iteration (Annex H); the revised version has not yet been tested.\\
\textit{Mechanic:} The first-iteration Rolling Out ritual required all participants to recite identical scripted phrases (``Thank you for listening to me,'' ``Thank you for speaking your truth''), producing a collective recitation effect (uniform, simultaneous, formulaic) incompatible with the personal and emergent character of the rest of the format.\\
\textit{Strength:} Strong. Multi-session; multi-report; behavioral disruption (mumbling, questioning designer); design response drafted but untested.

\end{enumerate}

\bigskip

%%% FIELD 2 %%%

\noindent\textit{Why Field 2 is a required characteristic.} Growth through play requires genuine intellectual and social challenge. The research question asks for a format that elevates transformative potential, but transformation through structured activity depends on productive pressure: the kind of challenge that stretches participants' capacities and motivates them to develop. The Inside Out format deliberately removes the win/loss evaluative structure that standard debate uses as its primary pressure mechanism, addressing the community hierarchy problem the research question identifies. But this design choice creates a central tension: does removing the competitive stakes also remove the cognitive and social demands that make debate a growth activity? Without productive pressure, repeated play (confirmed by Field 1) would be pleasantly novel but developmentally inert. Field 2 tests whether the pressure was redirected rather than removed, asking whether a non-competitive evaluative structure can sustain the challenge that transformative engagement requires.

\bigskip

\subsection*{Field 2: Productive Pressure Mechanisms}

\subsubsection*{Positive potential clues}

\begin{enumerate}

\item \textbf{Visibly high engagement and effort during speeches}\\
\textit{Points toward:} Cognitive challenge retained; participants investing effort proportional to competitive formats.\\
\textit{Relevance to field:} Sustained high effort in the absence of win/loss stakes is the most direct test of Field 2's central question: whether productive pressure has been successfully redirected. If participants invest comparable effort to standard competitive debate without a competitive outcome driving that investment, the format has demonstrated that productive pressure does not require a binary result.\\
\textit{Observed:} Yes (P1, P2, T2).\\
\textit{Evidence:} In the second playtest, all Value Defenders immediately adopted elaborate roleplay stances (second-person address, emotional tone, worldbuilding detail) suggesting high rhetorical investment (P2, T2). The first playtest also showed high investment: participants came up with elaborate and nuanced solutions, constructed careful arguments, and several delivered quite passionate speeches, fully ``try-hard'' play despite the unfamiliar format (P1).\\
\textit{Mechanic:} The format as a whole. Interaction mechanics in particular (the Huddle and Points of Information) were among the most liked and requested features across both playtests, suggesting they play an especially important role in sustaining engagement and effort.\\
\textit{Strength:} Strong. Behavioral.

\item \textbf{Strong performance quality: complex arguments, empathetic reasoning, strategic engagement}\\
\textit{Points toward:} Format generating the right kind of cognitive demand.\\
\textit{Relevance to field:} Performance quality differentiates productive pressure from mere social pressure: it tests whether the format's challenge demands complex thinking and empathetic reasoning, not just effort for its own sake. For Field 2, quality matters because growth requires challenge of the right kind: participants can try very hard at a low-complexity task without developing.\\
\textit{Observed:} Yes (P1, P2, D1, D2).\\
\textit{Evidence:} Participants described ``architecting solutions'' as ``very complex and interesting strategically'' (P1, D1). In the second playtest, participants navigated value hierarchies and inter-value tensions without prompting. The Subject's S0 and S2 speeches in the second playtest demonstrated sophisticated contextual reasoning (T2). Both sessions also produced strikingly high levels of acting and roleplay: players adopted character mannerisms, delivered emotionally engaged speeches, and one played a background soundtrack during their S1. A key question raised in debrief was whether this emotional register came at the expense of critical thinking, but the evidence suggests it did not: very elaborate solutions, practicalities-based argumentation, and careful value hierarchies were consistently present, indicating that critical thinking and emotional engagement coexisted rather than traded off (P1, P2, T2).\\
\textit{Mechanic:} The motion structure provides a specific character context that makes empathetic reasoning contextually available and strategically useful, rewarding perspective-taking as a competitive skill. Notably, the source of performance pressure shifts in this format: rather than the win/loss dynamic, the pressure becomes social: being recognized and appreciated by fellow Inner Voices and the Subject for the quality and resonance of one's contributions. This social pressure is intrinsically tied to the quality of the engagement rather than to a binary outcome, making it a more nuanced and arguably healthier motivator.\\
\textit{Strength:} Moderate. Partially self-report; partially observed; consistent across sessions.

\item \textbf{Extensive post-debate reflection without researcher prompting}\\
\textit{Points toward:} Format generating cognitive and emotional residue worth processing.\\
\textit{Relevance to field:} Unsolicited post-debate reflection indicates the format generated cognitive and emotional residue that participants feel compelled to process, the phenomenological signature of productive pressure that goes beyond task completion. For Field 2, this is evidence that the challenge architecture produces the kind of generative engagement participants continue working through after the formal activity ends.\\
\textit{Observed:} Yes (P1, P2, D1, D2).\\
\textit{Evidence:} Both debrief sessions extended well beyond structured questions. Participants in the first playtest returned to discussing the motion's dilemma long after the formal debrief. Participants in both sessions also discussed, with evident pride, how fun and interesting it had been to try out techniques that were new to them: the second-person address, the emotional appeals, the worldbuilding. This is a distinct form of reflection from dilemma-processing: it is reflection on craft and creative agency, suggesting the format activates not just cognitive and moral residue but performative pride in artistic experimentation.\\
\textit{Mechanic:} The whole format; in particular, the open-ended character of the roleplay task provides room for genuine creative experimentation that standard debate's tight genre conventions do not.\\
\textit{Strength:} Moderate. Observed; informally behavioral; consistent across both playtests.

\item \textbf{Expressed drive to improve; participants identifying their own growth edges}\\
\textit{Points toward:} Format creating the kind of productive pressure that motivates development.\\
\textit{Relevance to field:} Self-identified growth edges demonstrate that the format's pressure is experienced as intrinsically motivating rather than externally coercive: participants want to develop because the format has made them aware of capacities they do not yet have. For Field 2, this is the self-sustaining pressure indicator: a challenge architecture that makes participants their own most demanding evaluators has achieved the goal better than any external scoring system could.\\
\textit{Observed:} Yes (P1, D1, P2, D2).\\
\textit{Evidence:} Participants in the first session said they would aim to be ``more emotionally immersed'' on replay (P1, D1), identifying immersion as a dimension they had not yet mastered. In the second playtest, participants initiated a group conversation about a specific skill the format demands but they found difficult: being polite when doing counter-argumentation, navigating how to diminish the relevance of another Inner Voice's points without it feeling dismissive or adversarial (P2, D2). This reflects a more specific and format-native growth edge than emotional immersion, and was entirely self-identified. However, the very informal friends-hangout setting of the second session made the learning conversation less sustained and dynamic than it might have been in a more structured context; the topic was raised but not fully explored.\\
\textit{Mechanic:} The commendation system names specific skills (e.g., ``Believable Passion'') that participants can identify as aspirational targets; the roleplay dimension of the format provides a specific, named axis for development on replay that is recognized as distinct from standard debate skills. The polite counter-argumentation skill is format-specific: standard debate does not reward it, so recognizing it as a growth edge requires playing Inside Out.\\
\textit{Strength:} Moderate. Self-report; forward-looking; actual improvement behavior untested. The P2 polite counter-argumentation conversation is particularly interesting as a format-native, self-generated growth goal.

\end{enumerate}

\subsubsection*{Negative potential clues}

\begin{enumerate}

\item \textbf{Apathy or minimal effort during speeches}\\
\textit{Points away from:} Format generating meaningful productive pressure.\\
\textit{Relevance to field:} Apathy in a non-competitive context is the primary risk that the removal of win/loss stakes creates, and would constitute direct evidence that the design failed to redirect pressure successfully. For Field 2, the meaningful absence of apathy across both sessions is therefore strong supporting evidence: the motivational vacuum that removing competition could have created did not materialize.\\
\textit{Observed:} No.\\
\textit{Evidence:} Not manifested across either playtest. All participants engaged substantively with their speeches.\\
\textit{Strength:} N/A.

\item \textbf{Improvement discussion only when directly prompted by evaluator or researcher}\\
\textit{Points away from:} Intrinsic motivation to grow; would suggest format is experienced as externally evaluated rather than self-developing.\\
\textit{Relevance to field:} Growth discussion requiring external prompting would indicate that the format's pressure mechanism is experienced as externally imposed rather than internally generated. For Field 2, this would reveal a critical gap between the format producing challenge and the format producing self-directed growth orientation, the deeper claim transformative design requires.\\
\textit{Observed:} Partially (D1, D2).\\
\textit{Evidence:} Some debrief forms contained brief or minimal self-improvement reflections that appeared to be prompted rather than unprompted. However, other responses in both sessions were substantive and self-initiated. The pattern is inconclusive.\\
\textit{Strength:} Weak. Inconclusive; mixed within both debrief sets.

\item \textbf{Confusion about what to do within role; confusion about how to improve}\\
\textit{Points away from:} Immediate accessibility; clear pressure architecture.\\
\textit{Relevance to field:} Role confusion indicates the pressure architecture is opaque or inaccessible: participants may be motivated to improve but cannot identify how or toward what. For Field 2, this is a design-gap indicator that tests whether the pressure mechanism is scaffolded clearly enough to translate effort into directed growth, or whether it generates undirected effort that dissipates without producing development.\\
\textit{Observed:} Yes (P1, D1, P2).\\
\textit{Evidence:} In the first playtest, the Subject player described not knowing ``what to say'' or how to ``play very clean,'' uncertain how to hold both emotional investment and evaluative neutrality simultaneously (P1, D1). In the second playtest, the Subject ran out of prepared Huddle questions before the 10 minutes elapsed (P2), revealing that the role requires more scaffolding than the ruleset provided. Both gaps were addressed in the second design iteration.\\
\textit{Mechanic:} The S0 speech instructions in the first iteration did not specify the Subject's emotional register or clarify the relationship between roleplay inhabitation (appropriate in S0/S1) and evaluative synthesis (appropriate in S2); the Huddle preparation time did not scaffold the kinds of questions the Subject should prepare. Both gaps were directly addressed in the second design iteration's revised Subject role instructions.\\
\textit{Strength:} Strong as a design-gap indicator. Behavioral; multi-instance; directly theoretically informative.

\end{enumerate}

\bigskip

%%% FIELD 3 %%%

\noindent\textit{Why Field 3 is a required characteristic.} The research question targets specific, identifiable problems in existing debate practice (value convergence, information density, and competitive hierarchy) and demands a format whose growth outputs address those problems directly. It is insufficient for a format to be engaging (Field 1) and challenging (Field 2) if the challenge it produces develops the same narrow skill set as traditional formats: faster argument construction, adversarial rhetorical technique, and strategic point-scoring. The transformative potential the research question seeks to elevate is potential for genuine personal development: improved empathy, contextual reasoning, collaborative discourse, and perspective-taking. Field 3 tests whether the format's procedural rhetoric successfully encodes and transmits these skills through its rule structure, without explicit instruction.

\bigskip

\subsection*{Field 3: Growth Direction, Useful Skills and Perspectives}

\subsubsection*{Positive potential clues}

\begin{enumerate}

\item \textbf{Novel play patterns transferring real-world skills (empathy, contextual reasoning, negotiation)}\\
\textit{Points toward:} Format's procedural rhetoric successfully encoding and transmitting its intended skill set.\\
\textit{Relevance to field:} Uninstructed adoption of perspective-taking behaviors is the primary test of procedural rhetoric's effectiveness: if participants behave in the ways the format's procedural logic rewards without being told to, the rules are successfully encoding the desired skill orientation at the behavioral level. For Field 3, behavioral evidence of skill transfer is stronger than any self-report because it shows the format producing the relevant behaviors during play rather than generating retrospective claims about development.\\
\textit{Observed:} Yes (P1, P2, D1, D2, T2).\\
\textit{Evidence:} In the second playtest, all Value Defenders spontaneously adopted second-person address, direct eye contact with the Subject, and contextualizing worldbuilding details, none of which were instructed in the ruleset (P2, T2). In the first playtest, participants articulated learning to ``think about the context of decisions'' and to be ``more human and more honest about human nature'' (P1, D1). One participant noted the format develops ``conversational skills and empathy without necessarily having this educational background'' (P2, D2). Three additional novel play patterns are worth noting: (1) a consistently \textit{polite and collaborative form of counter-argumentation} emerged organically: players shifted focus rather than attacked, built on others rather than dismantling them, and recognized that outright rebuttal felt ``wrong'' in this context; (2) participants in the second playtest spontaneously moved toward \textit{alignment on a common solution} across the second round of speeches, a behavior directly inverse to standard debate and a paradigmatic instance of the format's procedural rhetoric encoding collaborative negotiation; (3) a \textit{push toward synthesis} was evident in both rounds, with players building on rather than contradicting prior contributions, mirroring the kind of constructive dialogue the format is designed to develop.\\
\textit{Mechanic:} The character assignment system (``you represent the Subject's value of X'') and the personal dilemma motion structure (providing a named character with specific circumstances) invite perspective-taking without instructing it; the Subject as a named individual rather than an abstract adjudicator provides a concrete human recipient for empathetic reasoning. The absence of a rebuttal-rewarding scoring category inverts the rhetorical incentive structure, making polite, collaborative, synthesis-oriented engagement strategically advantageous rather than weakly punished.\\
\textit{Strength:} Strong. Behavioral for second playtest; self-report for first playtest; consistent direction.

\item \textbf{Continuation of standard debate's positive skill patterns (research, structured argumentation)}\\
\textit{Points toward:} Format building on rather than displacing existing debate skills.\\
\textit{Relevance to field:} The preservation of existing valuable skills demonstrates the format builds on rather than displaces the intellectual infrastructure debate already develops well. For Field 3, this matters because the research question asks for a format that addresses value convergence without sacrificing the intellectual challenge of the activity: a format that develops empathy at the cost of rigorous argument construction would trade one problem for another.\\
\textit{Observed:} Yes (P1, P2).\\
\textit{Evidence:} Participants continued to employ structured argument construction and research-based reasoning within speeches. The format did not produce a drift toward purely emotional or anecdotal argument; the value-argumentation task retained intellectual rigor.\\
\textit{Mechanic:} The speech order (A1, B1, C1, S1, A2, B2, C2, S2) preserves the argument-sequencing logic of standard debate; the value-argumentation task still rewards research and logical structure, redirected toward value priority rather than policy consequence chains.\\
\textit{Strength:} Low. Observational; consistent across both playtests, but it may be due more to the profile of the participants as experienced debaters. 

\item \textbf{Good play (empathetic, collaborative, value-aware) visibly rewarded in Subject Speech, Oral Critique, and debrief}\\
\textit{Points toward:} Evaluative culture aligned with transformative design intent.\\
\textit{Relevance to field:} The evaluative culture's recognition of empathetic and value-aware play demonstrates that the format's feedback loop is aligned with its growth direction: the skills being developed are also the skills being rewarded. For Field 3, this is the systemic confirmation: not just that participants develop useful skills, but that the format's community names and incentivizes those skills, which is what makes growth direction durable rather than session-specific.\\
\textit{Observed:} Partially (P1, P2, D1, D2).\\
\textit{Evidence:} Evaluators acknowledged empathetic and value-aware play in both oral critiques. Both playtest oral critiques were very insightful and offered meaningful feedback that rewarded good play. However, in both playtests Subject decisions did more rather reward overly-analytical play, but that points more towards the profile of the participants as highly experienced debaters.\\
\textit{Mechanic:} The oral critique provides a structured feedback channel explicitly tied to the evaluator's observations; the commendation system names specific skills to reward. Both mechanisms are directionally correct but depend on evaluator familiarity with the system, which is immature at small scale.\\
\textit{Strength:} Moderate. Partially implemented; evaluative culture immature at small scale.

\end{enumerate}

\subsubsection*{Negative potential clues}

\begin{enumerate}

\item \textbf{Non-standard behaviors being rewarded that do not transfer real-world skills}\\
\textit{Points away from:} Skill development value; would suggest format incentivizes play that is idiosyncratic to Inside Out rather than transferable.\\
\textit{Relevance to field:} Rewarding behaviors idiosyncratic to Inside Out would indicate the format has created its own insular skill set: developmentally valuable only within the format's fictional frame and not transferable to the real-world contexts the research question is concerned with. For Field 3, this would be the central failure mode: a format confirming Fields 1, 2, and 4 while developing skills that only matter within itself.\\
\textit{Observed:} No.\\
\textit{Evidence:} Not manifested across either playtest. No evaluator praise for behaviors inconsistent with empathy, contextual reasoning, or collaborative discourse was recorded.\\
\textit{Strength:} N/A.

\item \textbf{Emergence of unhelpful new play patterns (e.g., hollow performative emotionality without substance)}\\
\textit{Points away from:} Skill transfer; would suggest format's emotional register invites performance rather than genuine reflection.\\
\textit{Relevance to field:} Emotional performance without substance would indicate the format's roleplay invitation is being taken as license for affect theater rather than genuine perspective inhabitation. For Field 3, this would reveal a gap between the format's intended procedural rhetoric (emotional engagement as a route to perspective-taking) and its actual output: emotional performance as a substitute for it.\\
\textit{Observed:} No.\\
\textit{Evidence:} Not manifested across either playtest. Emotional engagement observed in both sessions was substantively connected to argument construction and scenario reasoning.\\
\textit{Strength:} N/A.

\item \textbf{Immersion-breaking format elements preventing the bleed arc from completing}\\
\textit{Points away from:} Memetic and emancipatory bleed; a truncated immersion arc limits transformative potential.\\
\textit{Relevance to field:} The bleed arc is the mechanism through which in-game learning becomes available as a resource for real-world situations; if the transformational container breaks before bleed completes, the skills developed inside the format remain quarantined within the fictional frame rather than crossing the play/life boundary. For Field 3, this is the most structurally significant gap: in-round growth direction can be confirmed while the bleed arc failure prevents that growth from becoming genuinely transferable.\\
\textit{Observed:} Yes (P1, D1).\\
\textit{Evidence:} The ending ritual (Rolling Out) was experienced as ``very robotic'' and described as breaking immersion at the moment it was intended to complete the emotional arc (P1, D1). Because emancipatory bleed requires a sustained transformational container through closing, a ritual that breaks character at the end undermines the mechanism. This is the most significant growth-direction design gap identified.\\
\textit{Mechanic:} The first-iteration Rolling Out ritual required participants to recite identical scripted phrases in sequence, producing a uniform collective-recitation effect incompatible with the personal and emergent character of preceding play. The ritual was redesigned in the third design iteration (Annex H); that revised version has not yet been tested.\\
\textit{Strength:} Strong. Multi-report; behavioral disruption; theoretical mechanism clearly implicated; design response drafted but untested.

\end{enumerate}

\bigskip

%%% FIELD 4 %%%

\noindent\textit{Why Field 4 is a required characteristic.} The transformational container, the social space in which growth occurs, must be sustained not only within a single session but across the community that repeatedly inhabits the format. The research question identifies "community hierarchy" as a specific structural problem in existing formats, and transformative game design theory holds that the container's integrity is a prerequisite for bleed and sustained growth: participants will not take the vulnerable risks that genuine transformation requires if the surrounding community rewards status competition, punishes unconventional play, or makes failure socially costly. Fields 1--3 confirm that individual participants engage, are challenged, and develop useful skills within a session; Field 4 tests whether the format's structural logic reshapes the social environment those sessions occur within: whether it can convert an adversarial community culture into one that sustains rather than undermines the container.

\bigskip

\subsection*{Field 4: Community Dynamics and Reduced Toxicity}

\subsubsection*{Positive potential clues}

\begin{enumerate}

\item \textbf{Unsolicited expressions of safety and reduced anxiety}\\
\textit{Points toward:} Transformational container functioning; psychological safety achieved without sacrificing engagement.\\
\textit{Relevance to field:} Anxiety reduction reported without direct prompting, across independent groups, demonstrates that the format's structural modifications are producing the psychological safety effect, not researcher reassurance or social expectation. For Field 4, the source of safety matters: safety created by researcher presence would not persist in the format's independent use; safety created by structural logic is durable and community-portable.\\
\textit{Observed:} Yes (P1, P2, D1, D2).\\
\textit{Evidence:} ``Much much less anxiety-inducing than standard debate'' was consistently reported across both playtests and multiple independent participants (P1, P2, D1, D2). Participants described the experience as ``quieter, more rested'' (P1) and noted a ``feeling of salvation regardless of the result'' (P2, D2). The consistency across two independent groups strengthens confidence in this finding.\\
\textit{Mechanic:} Four mechanics operate jointly to reduce performance anxiety: (1) the Evaluator seating arrangement places the observer out-of-sight and peripheral, removing the primary anxiety-generating figure of standard debate; (2) no public scoring occurs during the round; (3) binary win/loss outcomes are removed, eliminating zero-sum social capital stakes; (4) the personal dilemma motion has no objectively correct answer, removing the possibility of being demonstrably wrong.\\
\textit{Strength:} Very strong. Multiple independent reports; consistent across two sessions; theoretically predicted prerequisite for bleed.

\item \textbf{Brave, vulnerable play acknowledged and rewarded}\\
\textit{Points toward:} Evaluative culture recognizing and incentivizing the emotional risk the format asks of participants.\\
\textit{Relevance to field:} Community acknowledgment of brave play demonstrates that the format's evaluative culture is aligning with its design intent in real time: unconventional, emotionally risky contributions are received as strengths rather than exploitable weaknesses. For Field 4, this is the confirmation that the container is actively maintained by the community rather than merely assumed by the format's rules.\\
\textit{Observed:} Yes (P1, D1, P2, D2).\\
\textit{Evidence:} Multiple instances of brave and unconventional play were observed across both playtests. In the first session, the C team chose a deliberately vulnerable and conversational approach (emotionally intimate and unhurried) that is highly uncommon in the standard formats those participants were trained in; this was noted by others and acknowledged in the oral critique (P1, D1). In the second session, brave play took several forms: the S1 player chose to play a background soundtrack during their opening speech, a choice with no precedent in debate and a real social risk; the C team built their entire case around a worldbuilt character (the Subject's mother) and her emotional stakes; players across the room adopted highly expressive, roleplay-heavy delivery styles; and participants responded to these choices with visible gesture-based approval and enthusiastic nonverbal acknowledgment throughout: laughing, nodding, leaning forward. The unplanned bleed event (a participant sharing privately that the scenario hit close to a real family situation) was also noted and acknowledged in the oral critique (P2, D2). Formal commendation assignment for brave play was not consistently enacted due to system immaturity.\\
\textit{Mechanic:} The bleed-check moment in the third-iteration Rolling Out ritual explicitly creates a designated space for emotional acknowledgment; the Evaluator role instructions frame psychological safety as a primary responsibility, making responsive acknowledgment of vulnerable play structurally expected. The removal of win/loss outcomes reduces the cost of unconventional choices, making brave play a lower-risk option.\\
\textit{Strength:} Strong; multi-instance; behavioral; evaluative response was informal rather than systematic.

\item \textbf{Participants articulating liberation from standard debate's performance pressure (emancipatory bleed)}\\
\textit{Points toward:} Emancipatory bleed occurring: participants experiencing the format as freeing rather than constraining, and building confidence in their ability to speak up.\\
\textit{Relevance to field:} Emancipatory bleed reports, participants experiencing the format as freeing from the evaluative regime of existing formats, demonstrate that the format is dismantling rather than merely replacing standard debate's performance anxiety architecture. For Field 4, emancipatory bleed is the mechanism that specifically addresses the "community hierarchy" problem: it is the experiential evidence that the format is reaching and liberating the voices that status competition suppresses.\\
\textit{Observed:} Yes (P1, P2, D1, D2).\\
\textit{Evidence:} ``More relaxed, not performing execrably, felt like playing really okay'' (P1, D1); ``gives a little feeling of salvation, regardless of the result'' (P2, D2); ``everyone distinguished himself by his own style, which in ordinary debate competitions does not happen'' (P2, D2). These reflections articulate a liberation from the evaluative regime of standard debate, congruent with previously reported clues of an overall much lower level of anxiety experienced by all participants.\\
\textit{Mechanic:} Character assignment externalizes personal value expression into a role (``I am arguing for the value of X''), reducing the social exposure of holding a minority position; the removal of binary outcomes eliminates the competitive cost of expressing an unconventional view; the personal style encouraged by the roleplay premise means individual expressiveness is rewarded rather than penalized.\\
\textit{Strength:} Strong. Consistent across both playtests; multiple independently-worded formulations of the same attitudinal shift.

\item \textbf{Emergent norms of collaborative rather than adversarial discourse}\\
\textit{Points toward:} Format's procedural rhetoric successfully reshaping community conduct without explicit instruction.\\
\textit{Relevance to field:} Spontaneous norm articulation (participants identifying without instruction that adversarial rebuttal "feels wrong") is the highest-strength behavioral indicator that the format's procedural rhetoric has reshaped social orientation at the community level, not just the individual level. For Field 4, emergent norms are the gold standard of structural success: they show the format's logic has been collectively internalized, producing community behavior the ruleset did not explicitly mandate.\\
\textit{Observed:} Yes (P2, D2).\\
\textit{Evidence:} Participants in the second playtest spontaneously articulated that ``outright attacking others' points feels wrong, disrespectful, and unnecessary.'' This was not a rule in the format but an emergent norm arising from the format's structure, a paradigmatic instance of procedural rhetoric in action. Players also noted that shifting focus and building on others' contributions felt ``more natural and more persuasive'' than traditional rebuttal.\\
\textit{Mechanic:} The Subject as shared persuasion target means all Value Defenders are competing for the same person's attention rather than defeating each other; no scoring category rewards rebuttal specifically, meaning attacking a competitor's argument carries no direct competitive benefit while building on it does, thus inverting the rhetorical incentive structure of standard competitive debate.\\
\textit{Strength:} Strong. Behavioral; emergent; unsolicited; theoretically predicted.

\end{enumerate}

\subsubsection*{Negative potential clues}

\begin{enumerate}

\item \textbf{Expressions of distress during or after the round}\\
\textit{Points away from:} Psychological safety; would suggest the format activates personal stakes without providing sufficient containment.\\
\textit{Relevance to field:} Distress in a format that deliberately elicits personal stakes would indicate the format lacks sufficient containment for the vulnerability it invites: the transformational container is failing its primary function. For Field 4, this is the most consequential failure mode: a format that elicits vulnerability and fails to hold it safely is not merely ineffective but actively harmful, which is why the absence of distress despite moments of genuine emotional intensity is particularly significant.\\
\textit{Observed:} No.\\
\textit{Evidence:} Not manifested across either playtest. The unplanned bleed event (participant sharing a real family situation) was handled calmly within the debrief structure; no distress was recorded. The bleed-check and debrief protocols appear to have functioned as intended.\\
\textit{Strength:} N/A. Notable absence given the emotional intensity of some moments.

\item \textbf{Arrogant, dismissive, or unempathetic behavior during debate}\\
\textit{Points away from:} Community health; toxicity reduction.\\
\textit{Relevance to field:} Adversarial social behavior from participants trained in competitive debate would indicate that the format's structural modifications have not successfully overridden the behavioral dispositions those participants bring from standard formats. For Field 4, this would reveal a gap between structural design and cultural practice: the format's rules may reward collaborative behavior while the community's habits continue to reward adversarial conduct.\\
\textit{Observed:} No.\\
\textit{Evidence:} Not manifested across either playtest. No dismissive or contemptuous behavior toward other participants was recorded in either field notes or transcripts.\\
\textit{Strength:} N/A.

\item \textbf{Such behavior being praised by the evaluator or community}\\
\textit{Points away from:} Evaluative culture aligned with design intent; would signal the format's recognition system is rewarding the wrong behaviors.\\
\textit{Relevance to field:} Praise of dismissive behavior would indicate the format's evaluative culture has been captured by existing norms rather than reshaping them: the recognition system, the primary lever for community change, is pointing in the wrong direction. For Field 4, the evaluative culture test is the most systemic indicator: if the community rewards the behaviors the format is designed to discourage, the social logic will prevail over the structural logic over time.\\
\textit{Observed:} No.\\
\textit{Evidence:} Not manifested across either playtest.\\
\textit{Strength:} N/A.

\end{enumerate}

\bigskip

%%% FIELD 5 %%%

\noindent\textit{Why Field 5 is a required characteristic.} The research question centers on "value convergence" as the primary epistemic problem the format must address, and epistemic humility (the capacity to hold beliefs as provisional, to recognize the legitimacy of competing value frameworks, and to treat disagreement as inquiry rather than contest) is its direct structural opposite. Fields 1--4 together demonstrate that the format is engaging, challenging, skill-building, and community-safe; but none of those fields tests whether the format addresses value convergence at the level of how participants hold their own beliefs. A format could confirm all four earlier fields while still producing participants who are enthusiastic, pressured, skilled, and socially safe but epistemically unchanged, still convinced that their values are objectively correct and others' are deficient. Field 5 tests the deepest targeted transformation: whether the format's procedural structure (Inner Voices framing, S0 bias disclosure, the absence of a winner, the multi-perspective structure) produces a measurable shift toward value pluralism as an emergent consequence of inhabiting the format's logic.

\bigskip

\subsection*{Field 5: Epistemic Humility}

\subsubsection*{Positive potential clues}

\begin{enumerate}

\item \textbf{Unprompted awareness that multiple competing perspectives hold genuine value}\\
\textit{Points toward:} Core attitudinal shift targeted by the format: epistemic humility arising from procedural engagement with value conflict.\\
\textit{Relevance to field:} Attitudinal shifts toward value plurality arising without researcher prompting demonstrate that the format's procedural structure produced the epistemic orientation the research question targets, rather than a researcher-advocated conclusion. For Field 5, the unprompted character of this evidence is essential: the research question asks whether the format can address value convergence structurally; if the shift requires prompting, it confirms only that the research team valued epistemic humility, not that the format produces it.\\
\textit{Observed:} Yes (P1, D1, P2, D2).\\
\textit{Evidence:} ``Every team says valuable stuff and it is important more to see what is important to us'' (P1, D1). This formulation, articulated without prompting after a single encounter, captures precisely the attitudinal shift the design targets: moving from ``which side is right'' to ``what do I value and why.'' Additional supporting evidence from the first playtest includes participants articulating that arguing someone is ``morally wrong'' for their personal decisions feels ``a bit absurd'' in this context, and recognition that the format makes personal value priorities legible rather than evaluating them against an external standard (P1, D1). In the second playtest, a notable extended conversation emerged in the debrief about how the Subject's opening speech (S0) institutionalises the practice of surfacing personal biases, making explicit what a player personally cares about before the debate even begins. Participants contrasted this with other formats, where the pretense of objectivity is maintained throughout: in those formats, they argued, biases are falsely supposed not to exist, propagating the idea that only one particular type of ``objective'' person can judge well and legitimizing social pressure to perform that identity. The S0 mechanism, by contrast, invites players to name their values openly from the start, reframing bias not as contamination but as perspective (P2, D2).\\
\textit{Mechanic:} The three-value motion structure assigns each team a value that is genuinely defensible for the character; the Subject's final synthesis requirement frames the question as ``what matters most to this person in this context'' rather than ``which position is objectively correct,'' making all three perspectives legitimate from the outset. The S0 speech institutionalises bias transparency rather than suppressing it, which makes epistemic humility structurally available rather than personally effortful.\\
\textit{Strength:} Moderate. Multi-session; self-report; directly on-target with design intent; P1 formulation notable for being unprompted; P2 debrief conversation particularly theoretically rich.

\item \textbf{Demonstrated understanding that personal values are context-dependent, not universal}\\
\textit{Points toward:} Internalization of the format's procedural rhetoric: that decisions are contextually embedded and legitimately vary by person and circumstance.\\
\textit{Relevance to field:} Recognizing that decisions are personal rather than universal is the specific intellectual content that distinguishes epistemic humility from mere tolerance: it is not just accepting that others hold different views but understanding that those views are legitimately contextual. For Field 5, this is the cognitive-structural indicator: the format has made contextual embeddedness of values experientially legible through play, not through assertion.\\
\textit{Observed:} Yes (P1, P2, D1, T2).\\
\textit{Evidence:} ``Sometimes decisions are personal, not universal'' (P1, D1). The Subject's S0 and S2 arc in the second playtest demonstrated a similar trajectory: moving from roleplay inhabitation (personal, contextual) toward evaluative synthesis (what values do I hold about this?), with explicit self-awareness of that transition (P2, T2).\\
\textit{Mechanic:} The character-specific motion context (personal history, relationships, circumstances) makes universal moral claims strategically inappropriate to the scenario; the S0 instructions ask the Subject to describe the scenario from a personal perspective before argumentation begins, embedding contextual reasoning before any abstract principle is introduced.\\
\textit{Strength:} Moderate. Multi-session; self-report and behavioral; theoretical coherence strong.

\item \textbf{Spontaneous movement toward synthesis rather than competitive winning}\\
\textit{Points toward:} Format orienting participants toward value integration rather than dominance: the deepest available behavioral indicator of epistemic humility in action.\\
\textit{Relevance to field:} Synthesis behavior emerging from the format's structural logic, without instruction and in inversion of standard debate's escalation pattern, is the deepest behavioral indicator that the epistemic shift has reached the level of action, not only self-perception. For Field 5, the behavioral character of this evidence is decisive: self-reports of attitudinal shift are susceptible to social desirability effects; transcript-confirmed spontaneous synthesis demonstrates the format has made integration rather than dominance the natural cognitive direction.\\
\textit{Observed:} Yes (P2, T2).\\
\textit{Evidence:} In the second round and also less in the first one of Value Defender speeches, participants naturally moved toward synthesis rather than intensifying conflict, the reverse of what standard debate produces. This was not instructed; it emerged from the format's procedural logic (P2, T2). Additionally, the Subject's opening statement ``I'm not sure if that's the solution; that's why this thing is eating at me'' (P2, T2, transcript-confirmed) demonstrates genuine epistemic uncertainty from the start of the round, produced by the format's demand that the Subject hold the dilemma rather than resolve it prematurely. Two further details reinforce this orientation toward collegiality and synthesis: (1) participants in the second playtest organically renamed themselves ``Inner Voices'' (which was subsequently formalized in the third design iteration), a label that positions them as parts of a shared internal deliberation rather than as opposing sides, colleagues rather than opponents; (2) the rules themselves permit multiple Inner Voices to argue for the same solution, structurally allowing consensus rather than requiring disagreement, which removes the adversarial constraint that standard debate formats impose.\\
\textit{Mechanic:} Speeches A2/B2/C2 are expected to build on A1/B1/C1 rather than dismantle competitors' positions. The Subject's S1 is an explicitly preliminary decision requiring ongoing engagement, structurally rewarding continued dialogue. The Huddle allows the Subject to ask clarifying questions, modeling collaborative inquiry as a legitimate competitive move. The Inner Voices naming and the same-solution permission together reframe competition from adversarial to collegial at the level of both identity and rules.\\
\textit{Strength:} Strong. Behavioral; transcript-confirmed; emergent; theoretically predicted.

\end{enumerate}

\subsubsection*{Negative potential clues}

\begin{enumerate}

\item \textbf{Unequivocal condemnation of beliefs or belief-holders}\\
\textit{Points away from:} Epistemic humility; would indicate the format is activating moral certainty rather than productive uncertainty.\\
\textit{Relevance to field:} Condemnation of positions or belief-holders would demonstrate that the format activates moral certainty rather than productive uncertainty: instead of making epistemic humility the natural cognitive orientation, the format's emotional intensity is channeling into absolutism. For Field 5, this is the mirror-image failure mode: a format that elicits strong value engagement could either produce humility (the target) or produce entrenchment (its opposite), and this clue tests which outcome the format's structural logic generates.\\
\textit{Observed:} No.\\
\textit{Evidence:} Not manifested across either playtest or in the transcripts. No unequivocal condemnation of a position, values system, or person was recorded.\\
\textit{Strength:} N/A. Absence is notable given that the motions used both involved choices with genuine ethical weight.

\item \textbf{Such condemnation being normalized or praised}\\
\textit{Points away from:} Community epistemic culture; would signal the format's evaluative frame rewards dogmatic certainty.\\
\textit{Relevance to field:} Normalization of condemnation would indicate the epistemic failure has reached the community level: not just that individuals experience certainty but that the evaluative culture rewards it. For Field 5, the evaluative culture test is the systemic indicator: the research question's goal of a format that "addresses the issues of value convergence" requires the format's community to sustain epistemic humility collectively, not only momentarily within individuals.\\
\textit{Observed:} No.\\
\textit{Evidence:} Not manifested across either playtest.\\
\textit{Strength:} N/A.

\end{enumerate}
